%----------------------------------------------------------------------------------------
% TOEPASSINGEN VAN AFGELEIDEN. Alle opgaven zijn eigen werk.
%
% Opbouw: Start, leerstof (theorie, werkwijze of uitleg, voorbeeld, korte oefening), en op een
% nieuwe pagina de Opdrachten (door elkaar, stijgende moeilijkheid), Koppel en Test jezelf.
% Labels van dit hoofdstuk beginnen met h4-.
%----------------------------------------------------------------------------------------

\mychapter{4}{Toepassingen van afgeleiden}

\lkr{Leerplan Basis: LPD 14-15 (afgeleiden), het deel over het verloop van functies (stijgen, dalen, extrema, holheid, buigpunten) en optimaliseringsproblemen. Over \textit{hol} en \textit{bol} bestaan verschillende afspraken; dit boek koppelt ze telkens aan de vorm van een dal- of bergparabool, zodat er geen misverstand kan zijn. De stelling dat $f' > 0$ stijgen geeft, wordt hier niet bewezen (ze volgt uit de middelwaardestelling).}

In Algemene functies las je toppen en dalen af op een grafiek. Nu kun je ze berekenen: in een top of een dal is de raaklijn horizontaal, dus de afgeleide nul. Daarmee los je ook praktische vragen op: hoe maak je een blik met zo weinig mogelijk metaal, of een weide zo groot mogelijk?

\begin{oefeningen}[Start]
\oef{h4-start-1}

\oef{h4-verken}
\end{oefeningen}

%========================================================================================
\section{Stijgen, dalen en extrema}

\subsection{Stijgen en dalen}

Een functie $f$ is stijgend op een interval $I$ als een grotere $x$ altijd een grotere functiewaarde geeft: uit $x_1 < x_2$ volgt $f(x_1) < f(x_2)$. Ze is dalend als uit $x_1 < x_2$ volgt dat $f(x_1) > f(x_2)$.

\naast{%
Waar een grafiek stijgt, lopen de raaklijnen omhoog: hun richtingscoëfficiënt is positief. Waar ze daalt, lopen ze omlaag. Het teken van de afgeleide vertelt dus of de functie stijgt of daalt. Dat is het idee achter de volgende stelling.}{%
\begin{tikzpicture}[x=.7cm,y=.45cm,font=\scriptsize]
  \draw[->,grijs] (-.3,0) -- (5.4,0) node[right,text=black] {$x$};
  \draw[->,grijs] (0,-.3) -- (0,4.4) node[above,text=black] {$y$};
  \kromme{.2}{5}{2+1.5*sin(70*(\x-.3))}
  \draw[grijs] (.2,1.9) -- (1.2,3.51);
  \draw[grijs] (3.1,1.41) -- (4.1,.26);
  \node[grijs] at (1.5,1.5) {$f' > 0$};
  \node[grijs] at (4.6,2.4) {$f' < 0$};
\end{tikzpicture}}

\begin{theorie}
\begin{stelling}{Teken van de afgeleide}
Is $f$ afleidbaar op een interval $I$, dan geldt: \\[1mm]
\qquad is $\formule{f'(x) > 0}$ voor alle $x \in I$, dan is $f$ \textbf{\textcolor{accent}{stijgend}} op $I$; \\
\qquad is $\formule{f'(x) < 0}$ voor alle $x \in I$, dan is $f$ \textbf{\textcolor{accent}{dalend}} op $I$.
\end{stelling}
\end{theorie}

Dat lijkt vanzelfsprekend, maar het bewijs is niet eenvoudig. Het steunt op de middelwaardestelling van Lagrange, die buiten deze cursus valt. We nemen de stelling aan.

Let op: de stelling werkt maar in één richting. De functie $x^3$ stijgt overal, maar haar afgeleide $3x^2$ is nul in $0$. Een positieve afgeleide is dus voldoende om te stijgen, maar niet nodig.

\voorbeeld{}{Waar stijgt en waar daalt $f(x) = x^2 - 4x + 1$?\label{vb:h4-stijgen}}{%
\naast{%
\uitleg{Bepaal het teken van de afgeleide.}
\begin{stappenvrij}{3.9cm}{5.4cm}
f'(x) = 2x - 4 & De afgeleide, een eerstegraadsfunctie. \\
f'(x) < 0 \iff x < 2 & Daar daalt $f$: op $\left]-\infty, 2\right[$. \\
f'(x) > 0 \iff x > 2 & Daar stijgt $f$: op $\left]2, +\infty\right[$.
\end{stappenvrij}
Dat klopt met de grafiek: een dalparabool met top in $x = 2$.}{%
\begin{tikzpicture}[x=.5cm,y=.35cm,font=\scriptsize]
  \rooster{-1}{1}{5}{-3}{1}{4}
  \kromme{-.6}{4.6}{(\x)^2-4*\x+1}
  \vol{(2,-3)}
\end{tikzpicture}}}

\begin{oefeningen}
\oef{h4-stijgen}
\end{oefeningen}

\subsection{Extrema}

In een top of een dal van een gladde grafiek is de raaklijn horizontaal. Dat zag je in de Verken. Eerst leggen we vast wat een top en een dal precies zijn: het hoogste of laagste punt in de buurt, niet noodzakelijk van de hele grafiek.

\begin{theorie}
$f(a)$ is een \textbf{\textcolor{accent}{maximum}} van $f$ als $\formule{f(x) \leqslant f(a)}$ voor alle $x$ in een interval rond $a$, en een \textbf{\textcolor{accent}{minimum}} als $\formule{f(x) \geqslant f(a)}$ voor alle $x$ in een interval rond $a$. Een maximum of minimum heet een \textbf{\textcolor{accent}{extremum}}.

\begin{stelling}{Extremum en afgeleide}
Heeft $f$ een extremum in $a$, en is $f$ afleidbaar in $a$, dan is $\formule{f'(a) = 0}$.
\end{stelling}
\begin{bewijs}
Stel dat $f(a)$ een maximum is. Voor kleine $h \neq 0$ is dan $f(a + h) - f(a) \leqslant 0$. Kijk naar het teken van het differentiequotiënt:
\bewijsregel{h > 0\colon\ \frac{f(a + h) - f(a)}{h} \leqslant 0 \qquad\qquad h < 0\colon\ \frac{f(a + h) - f(a)}{h} \geqslant 0 .}
Een limiet van getallen die allemaal $\leqslant 0$ zijn, kan niet positief zijn: anders zouden de getallen dicht bij die limiet ook positief zijn. De rechterlimiet van het differentiequotiënt is dus $\leqslant 0$, en net zo is de linkerlimiet $\geqslant 0$. Beide zijn gelijk aan $f'(a)$, dus $f'(a) \leqslant 0$ en $f'(a) \geqslant 0$, en $f'(a) = 0$. Is $f(a)$ een minimum, pas dit dan toe op $-f$, die in $a$ een maximum heeft.
\end{bewijs}
\end{theorie}

De extrema van een afleidbare functie vind je dus tussen de nulwaarden van $f'$. Maar niet elke nulwaarde is een extremum, en niet elk extremum geeft een nulwaarde:

\figuren{\begin{tikzpicture}[x=.5cm,y=.5cm,font=\scriptsize]
  \roosterleeg{-2}{1}{2}{-2}{1}{2}
  \kromme{-1.3}{1.3}{(\x)^3}
  \draw[grijs] (-1,0) -- (1,0);
  \node[below] at (0,-2.1) {$x^3$};
\end{tikzpicture}\qquad\qquad
\begin{tikzpicture}[x=.5cm,y=.5cm,font=\scriptsize]
  \roosterleeg{-2}{1}{2}{-2}{1}{2}
  \draw[accent,thick] (-2,2) -- (0,0) -- (2,2);
  \vol{(0,0)}
  \node[below] at (0,-2.1) {$|x|$};
\end{tikzpicture}\qquad\qquad
\begin{tikzpicture}[x=.5cm,y=.5cm,font=\scriptsize]
  \roosterleeg{-1}{1}{3}{-1}{1}{3}
  \kromme{0}{2}{\x+.25*(\x)^2}
  \vol{(0,0)} \vol{(2,3)}
  \node[below] at (1,-1.1) {$x + \frac{x^2}{4}$ op $[0, 2]$};
\end{tikzpicture}}

Bij $x^3$ is $f'(0) = 0$, maar de functie stijgt door: geen extremum. Bij $|x|$ is $0$ een minimum, maar daar is $f$ niet afleidbaar. En op een gesloten interval kan het grootste punt aan de rand liggen, waar de grafiek gewoon doorloopt. De stelling zegt alleen: \textit{als} er een extremum is waar $f$ afleidbaar is, dan is de afgeleide daar nul.

\subsection{Maximum of minimum?}

Of een nulwaarde van $f'$ een maximum, een minimum of geen van beide is, zie je aan het teken van $f'$ links en rechts ervan.

\begin{theorie}
\begin{stelling}{Tekenwissel van de afgeleide}
Is $f$ continu in $a$ en afleidbaar in een interval rond $a$ (behalve misschien in $a$), dan geldt: \\[1mm]
\qquad wisselt $f'$ in $a$ van $\formule{+}$ naar $\formule{-}$, dan is $f(a)$ een maximum; \\
\qquad wisselt $f'$ in $a$ van $\formule{-}$ naar $\formule{+}$, dan is $f(a)$ een minimum.
\end{stelling}
\begin{bewijs}
Stel dat $f' > 0$ op $\left]a - \delta, a\right[$ en $f' < 0$ op $\left]a, a + \delta\right[$. Neem $x$ links van $a$. Op $\left]a - \delta, a\right[$ stijgt $f$, dus voor $x < t < a$ is $f(x) < f(t)$. Laat $t$ naar $a$ gaan: omdat $f$ continu is in $a$, volgt
\bewijsregel{f(x) \leqslant f(a) .}
Neem $x$ rechts van $a$. Op $\left]a, a + \delta\right[$ daalt $f$, dus voor $a < t < x$ is $f(t) > f(x)$, en $t \to a$ geeft weer $f(a) \geqslant f(x)$. Dus $f(a)$ is een maximum. Bij een wissel van $-$ naar $+$ pas je dit toe op $-f$.
\end{bewijs}
\end{theorie}

Wisselt $f'$ niet van teken, dan is er geen extremum: de functie stijgt aan beide kanten, of daalt aan beide kanten. Alles samen vind je de extrema met een tekentabel van $f'$.

\voorbeeld{}{Bepaal het verloop en de extrema van $f(x) = x^3 - 3x$.\label{vb:h4-extrema}}{%
\naast{%
\uitleg{Leid af, ontbind, en stel de tekentabel van $f'$ op. Zet eronder pijlen voor $f$.}
\begin{stappenvrij}{4.2cm}{4.4cm}
f'(x) = 3x^2 - 3 & De machtregel. \\
= 3(x - 1)(x + 1) & Nulwaarden $-1$ en $1$. \\
\text{max. } f(-1) = 2 & $f'$: van $+$ naar $-$. \\
\text{min. } f(1) = -2 & $f'$: van $-$ naar $+$.
\end{stappenvrij}
\figuren{\tekentabel{cccccc}{x & & -1 & & 1 & \\\hline f'(x) & + & 0 & - & 0 & +}{f(x) & \nearrow & 2 & \searrow & -2 & \nearrow}}}{%
\begin{tikzpicture}[x=.6cm,y=.42cm,font=\scriptsize]
  \rooster{-2}{1}{2}{-3}{1}{3}
  \kromme{-2.1}{2.1}{(\x)^3-3*\x}
  \vol{(-1,2)} \vol{(1,-2)}
\end{tikzpicture}}}

\begin{werkwijze}
Bepaal het domein en leid af. Zoek de nulwaarden van $f'$ (en de punten waar $f'$ niet bestaat). Stel een tekentabel op voor $f'$, met daaronder pijlen voor $f$. Een wissel van $+$ naar $-$ is een maximum, van $-$ naar $+$ een minimum. Bereken de functiewaarden.
\end{werkwijze}

Een nulwaarde van $f'$ zonder tekenwissel kom je vaak tegen bij een kwadraat in de afgeleide.

\voorbeeld{}{Heeft $f(x) = x^3 - 3x^2 + 3x + 1$ extrema?\label{vb:h4-geenextr}}{%
\naast{%
\uitleg{Zoek de nulwaarden van $f'$ en kijk of het teken wisselt.}
\begin{stappenvrij}{4.2cm}{5cm}
f'(x) = 3x^2 - 6x + 3 & Afleiden. \\
= 3(x - 1)^2 & Een volkomen kwadraat. \\
f'(x) \geqslant 0 & Positief, behalve in $1$: daar nul, zonder tekenwissel.
\end{stappenvrij}
Geen extrema: $f$ stijgt overal. In $(1, 2)$ is de raaklijn horizontaal, maar de grafiek gaat er gewoon door, zoals bij $x^3$.}{%
\begin{tikzpicture}[x=.6cm,y=.42cm,font=\scriptsize]
  \rooster{-1}{1}{3}{-2}{1}{4}
  \kromme{-.55}{2.55}{(\x)^3-3*(\x)^2+3*\x+1}
  \draw[grijs] (.2,2) -- (1.8,2);
  \vol{(1,2)}
\end{tikzpicture}}}

\begin{oefeningen}
\oef{h4-tekentabel}
\end{oefeningen}

Met de afgeleiden uit hoofdstuk 3 kun je ook functies met $\e^x$, $\ln x$ of $\sin x$ onderzoeken. De tekentabel wordt dan vaak eenvoudiger dan je denkt, omdat sommige factoren altijd positief zijn.

\voorbeeld{}{Bepaal de extrema van $f(x) = x\,\e^{-x}$.\label{vb:h4-extrema2}}{%
\naast{%
\uitleg{Product- en kettingregel. De factor $\e^{-x}$ is altijd positief: het teken hangt alleen af van $1 - x$.}
\begin{stappenvrij}{4.5cm}{4.5cm}
f'(x) = \e^{-x} - x\,\e^{-x} & Productregel, met $(\e^{-x})' = -\e^{-x}$. \\
= (1 - x)\,\e^{-x} & $\e^{-x}$ buiten haakjes. \\
f' > 0 \text{ voor } x < 1 & Daar stijgt $f$. \\
f' < 0 \text{ voor } x > 1 & Daar daalt $f$. \\
\text{max. } f(1) = \tfrac{1}{\e} \approx 0{,}37 & Het enige extremum.
\end{stappenvrij}}{%
\begin{tikzpicture}[x=.6cm,y=2cm,font=\scriptsize]
  \draw[black!14,very thin,xstep=1,ystep=.25] (-1,-.5) grid (5,.5);
  \draw[->,grijs] (-1,0) -- (5.4,0) node[right,text=black] {$x$};
  \draw[->,grijs] (0,-.5) -- (0,.62) node[above,text=black] {$y$};
  \foreach \t in {1,2,3,4}{\node[below,inner sep=1.5pt] at (\t,0) {$\t$};}
  \node[left,inner sep=1.5pt] at (0,.5) {$0{,}5$};
  \node[left,inner sep=1.5pt] at (0,-.5) {$-0{,}5$};
  \kromme{-.4}{5}{\x*exp(-\x)}
  \vol{(1,.368)}
\end{tikzpicture}}}

\begin{oefeningen}
\oef{h4-extrema}

\oef{h4-extremaexp}
\end{oefeningen}

\subsection{Grootste en kleinste waarde op een interval}

Op een gesloten interval $[a, b]$ heeft een continue functie altijd een grootste en een kleinste waarde. Dat bewijzen we niet: net als bij de tussenwaardestelling steunt het op de volledigheid van de reële getallen. Die waarden liggen ofwel in een extremum binnen het interval, ofwel aan een rand.

\voorbeeld{}{Bepaal de grootste en de kleinste waarde van $f(x) = x^3 - 6x^2 + 9x$ op $[0, 5]$.\label{vb:h4-interval}}{%
\naast{%
\uitleg{De kandidaten zijn de nulwaarden van $f'$ in het interval en de twee randen. Vergelijk hun functiewaarden.}
\begin{stappenvrij}{4.2cm}{5cm}
f'(x) = 3(x - 1)(x - 3) & Nulwaarden $1$ en $3$, allebei in $[0, 5]$. \\
f(0) = 0, \ f(1) = 4 & De linkerrand en het maximum. \\
f(3) = 0, \ f(5) = 20 & Het minimum en de rechterrand. \\
\text{grootste: } 20 & In $x = 5$, aan de rand. \\
\text{kleinste: } 0 & In $x = 0$ en in $x = 3$.
\end{stappenvrij}
Het maximum $f(1) = 4$ is dus niet de grootste waarde op het interval.}{%
\begin{tikzpicture}[x=.55cm,y=.1cm,font=\scriptsize]
  \draw[black!14,very thin,xstep=1,ystep=5] (0,0) grid (5,20);
  \draw[->,grijs] (0,0) -- (5.4,0) node[right,text=black] {$x$};
  \draw[->,grijs] (0,0) -- (0,22) node[above,text=black] {$y$};
  \foreach \t in {1,2,3,4,5}{\node[below,inner sep=1.5pt] at (\t,0) {$\t$};}
  \foreach \t in {5,10,15,20}{\node[left,inner sep=1.5pt] at (0,\t) {$\t$};}
  \kromme{0}{5}{(\x)^3-6*(\x)^2+9*\x}
  \vol{(0,0)} \vol{(1,4)} \vol{(3,0)} \vol{(5,20)}
\end{tikzpicture}}}

\begin{werkwijze}
Zoek de nulwaarden van $f'$ die in het interval liggen (en de punten waar $f'$ niet bestaat). Bereken de functiewaarden in die punten en in de twee randen. De grootste en de kleinste daarvan zijn de gevraagde waarden.
\end{werkwijze}

\begin{oefeningen}
\oef{h4-interval}
\end{oefeningen}

%========================================================================================
\section{Holheid en buigpunten}

\subsection{De tweede afgeleide}

De afgeleide $f'$ is zelf een functie, en die kun je opnieuw afleiden. Zo ontstaat de tweede afgeleide. Ze zegt hoe snel de helling van $f$ verandert.

\begin{theorie}
De \textbf{\textcolor{accent}{tweede afgeleide}} van $f$ is de afgeleide van $f'$: \quad $\formule{f'' = (f')'}$. \\[1mm]
In de notatie van Leibniz, met $y = f(x)$: \quad $\formule{f''(x) = \frac{\mathrm{d}^2 y}{\mathrm{d}x^2}}$.
\end{theorie}

In een beweging is dat een bekende grootheid. Is $s(t)$ de plaats, dan is $v(t) = s'(t)$ de snelheid, en $a(t) = v'(t) = s''(t)$ de \textit{versnelling}: hoe snel de snelheid verandert.

\voorbeeld{}{Een lift vertrekt en stopt na $4$ s. Zijn hoogte is $s(t) = 6t^2 - t^3$ (in m, $0 \leqslant t \leqslant 4$). Bereken de snelheid en de versnelling. Wanneer rijdt hij het snelst?\label{vb:h4-tweede}}{%
\uitleg{Leid twee keer af. De snelheid is maximaal waar haar afgeleide, de versnelling, van teken wisselt.}
\begin{stappen}
v(t) = s'(t) = 12t - 3t^2 & De snelheid. Ze is $0$ bij het vertrek en in $t = 4$. \\[1ex]
a(t) = s''(t) = 12 - 6t & De versnelling. \\[1ex]
a(t) > 0 \text{ voor } t < 2, \ a(t) < 0 \text{ voor } t > 2 & Eerst versnelt de lift, daarna remt hij af. \\[1ex]
v(2) = 24 - 12 = 12 & De versnelling wisselt van $+$ naar $-$: de snelheid is maximaal na $2$ s, $12$ m/s.
\end{stappen}}

\begin{oefeningen}
\oef{h4-tweede}
\end{oefeningen}

\subsection{Hol en bol}

\naast{%
Bekijk de raaklijnen aan een dalparabool, van links naar rechts. Hun richtingscoëfficiënt neemt toe: $f'$ stijgt, en de grafiek buigt naar boven. Bij een bergparabool neemt ze af: $f'$ daalt, en de grafiek buigt naar onder. Dat gebruiken we als definitie.}{%
\begin{tikzpicture}[x=.5cm,y=.5cm,font=\scriptsize]
  \kromme{-2}{2}{.5*(\x)^2}
  \foreach \a in {-1.5,0,1.5}{\draw[grijs] ({\a-.7},{.5*\a*\a-.7*\a}) -- ({\a+.7},{.5*\a*\a+.7*\a});}
  \node[below] at (0,-.4) {hol};
\end{tikzpicture}\qquad
\begin{tikzpicture}[x=.5cm,y=.5cm,font=\scriptsize]
  \kromme{-2}{2}{2-.5*(\x)^2}
  \foreach \a in {-1.5,0,1.5}{\draw[grijs] ({\a-.7},{2-.5*\a*\a+.7*\a}) -- ({\a+.7},{2-.5*\a*\a-.7*\a});}
  \node[below] at (0,-.4) {bol};
\end{tikzpicture}}

\begin{theorie}
De grafiek van $f$ is \textbf{\textcolor{accent}{hol}} op een interval $I$ als $f'$ stijgend is op $I$: ze buigt naar boven, zoals een dalparabool ($\cup$). Ze is \textbf{\textcolor{accent}{bol}} op $I$ als $f'$ dalend is: ze buigt naar onder, zoals een bergparabool ($\cap$).

\begin{stelling}{Teken van de tweede afgeleide}
Is $f$ twee keer afleidbaar op een interval $I$, dan geldt: \\[1mm]
\qquad is $\formule{f''(x) > 0}$ voor alle $x \in I$, dan is de grafiek hol op $I$; \\
\qquad is $\formule{f''(x) < 0}$ voor alle $x \in I$, dan is de grafiek bol op $I$.
\end{stelling}
\begin{bewijs}
Pas de stelling over het teken van de afgeleide toe op de functie $f'$, met afgeleide $f''$. Is $f'' > 0$ op $I$, dan stijgt $f'$ op $I$, en dat is de definitie van hol. Is $f'' < 0$, dan daalt $f'$: bol.
\end{bewijs}
\end{theorie}

Zoals $f'$ zegt of de grafiek stijgt of daalt, zegt $f''$ dus hoe ze buigt. Over de woorden \textit{hol} en \textit{bol} bestaan verschillende afspraken: onthoud ze daarom met de figuur, $\cup$ en $\cap$.

\voorbeeld{}{Waar is de grafiek van $f(x) = x^3 - 3x$ hol, en waar bol?\label{vb:h4-holbol}}{%
\naast{%
\uitleg{Leid twee keer af en bepaal het teken van $f''$.}
\begin{stappenvrij}{3.9cm}{5.2cm}
f'(x) = 3x^2 - 3 & Eerste afgeleide. \\
f''(x) = 6x & Tweede afgeleide. \\
f''(x) < 0 \text{ voor } x < 0 & Bol op $\left]-\infty, 0\right[$: daar ligt de top. \\
f''(x) > 0 \text{ voor } x > 0 & Hol op $\left]0, +\infty\right[$: daar ligt het dal.
\end{stappenvrij}}{%
\begin{tikzpicture}[x=.55cm,y=.4cm,font=\scriptsize]
  \rooster{-2}{1}{2}{-3}{1}{3}
  \kromme{-2.1}{2.1}{(\x)^3-3*\x}
  \vol{(0,0)}
\end{tikzpicture}}}

\begin{oefeningen}
\oef{h4-holbol}
\end{oefeningen}

\subsection{Buigpunten}

In het vorige voorbeeld gaat de grafiek in de oorsprong over van bol naar hol. Zo'n punt heeft een eigen naam.

\begin{theorie}
Een \textbf{\textcolor{accent}{buigpunt}} van de grafiek van $f$ is een punt $(a, f(a))$ waar $f''$ van teken wisselt: de grafiek gaat er over van hol naar bol, of omgekeerd.

\begin{stelling}{Buigpunt en tweede afgeleide}
Is $f$ twee keer afleidbaar en is $(a, f(a))$ een buigpunt, dan is $\formule{f''(a) = 0}$.
\end{stelling}
\begin{bewijs}
Stel dat $f''$ in $a$ van $+$ naar $-$ wisselt. De functie $f'$ is afleidbaar, dus continu, en haar afgeleide $f''$ wisselt in $a$ van $+$ naar $-$. Volgens de stelling over de tekenwissel heeft $f'$ dus een maximum in $a$. Volgens de stelling over extrema is dan $(f')'(a) = f''(a) = 0$. Wisselt $f''$ van $-$ naar $+$, dan heeft $f'$ een minimum in $a$, met hetzelfde besluit.
\end{bewijs}
\end{theorie}

Een buigpunt is dus een punt waar de helling $f'$ een extremum heeft: daar is de grafiek het steilst of het minst steil van de buurt. Het omgekeerde geldt niet: $f(x) = x^4$ heeft $f''(x) = 12x^2$, nul in $0$, maar $f''$ wisselt er niet van teken. De grafiek is overal hol, en $(0, 0)$ is geen buigpunt.

\voorbeeld{}{Bepaal de extrema, de holheid en de buigpunten van $f(x) = x^4 - 4x^3$.\label{vb:h4-buigpunt}}{%
\naast{%
\uitleg{Stel twee tekentabellen op: voor $f'$ en voor $f''$. In $0$ is $f'$ nul zonder van teken te wisselen: daar is geen extremum.}
\begin{stappenvrij}{4.3cm}{3.2cm}
f'(x) = 4x^3 - 12x^2 & \\
= 4x^2(x - 3) & Nulwaarden $0$ en $3$. \\
\text{min. } f(3) = -27 & Geen extremum in $0$. \\
f''(x) = 12x^2 - 24x & \\
= 12x(x - 2) & Nulwaarden $0$ en $2$. \\
\text{buigpunten } (0, 0),\ (2, -16) & $f''$ wisselt van teken.
\end{stappenvrij}
\figuren{\tekentabel{cccccccc}{x & & 0 & & 2 & & 3 & \\\hline f'(x) & - & 0 & - & - & - & 0 & + \\ f''(x) & + & 0 & - & 0 & + & + & +}{f(x) & \cup & 0 & \cap & -16 & \cup & -27 & \cup}}}{%
\begin{tikzpicture}[x=.55cm,y=.075cm,font=\scriptsize]
  \draw[black!14,very thin,xstep=1,ystep=10] (-1,-30) grid (4,20);
  \draw[->,grijs] (-1.3,0) -- (4.4,0) node[right,text=black] {$x$};
  \draw[->,grijs] (0,-31) -- (0,22) node[above,text=black] {$y$};
  \foreach \t in {-1,1,2,3,4}{\node[below,inner sep=1.5pt] at (\t,0) {$\t$};}
  \foreach \t in {-20,-10,10}{\node[left,inner sep=1.5pt] at (0,\t) {$\t$};}
  \kromme{-1.2}{4.25}{(\x)^4-4*(\x)^3}
  \vol{(3,-27)} \vol{(0,0)} \vol{(2,-16)}
\end{tikzpicture}}}

\begin{werkwijze}
Bereken $f''$ en zoek de nulwaarden. Stel een tekentabel op voor $f''$: waar $f'' > 0$ is de grafiek hol, waar $f'' < 0$ bol. Een nulwaarde waar $f''$ van teken wisselt, geeft een buigpunt; bereken de functiewaarde.
\end{werkwijze}

\begin{oefeningen}
\oef{h4-buigpunt}

\oef{h4-buigpunt2}
\end{oefeningen}

\section{Optimaliseren}

Bij een optimaliseringsprobleem zoek je de beste keuze: de grootste oppervlakte, het minste materiaal, de kortste tijd. Wiskundig is dat telkens hetzelfde: de grootste of kleinste waarde van een functie. Het afleiden is meestal het makkelijkste deel. Het werk zit in het opstellen van die functie, met één veranderlijke, en in het bepalen van het domein.

\voorbeeld{}{Een boer zet een rechthoekige weide af langs een rivier. Aan de rivierkant komt geen hek; voor de drie andere zijden heeft hij $100$ m draad. Hoe groot kan de weide hoogstens worden?\label{vb:h4-weide}}{%
\naast{%
\uitleg{Er zijn twee onbekende zijden, maar de draad legt een verband op. Kies de breedte $x$ loodrecht op de rivier. Dan blijft er voor de lengte langs de rivier $100 - 2x$ over.}
\begin{stappenvrij}{4.6cm}{4.6cm}
A(x) = x(100 - 2x) & De oppervlakte, als functie van één veranderlijke. \\
0 < x < 50 & Het domein: beide zijden moeten positief zijn. \\
A'(x) = 100 - 4x & Afleiden. \\
A'(x) = 0 \iff x = 25 & De enige kandidaat. \\
A' > 0 \text{ voor } x < 25 & \\
A' < 0 \text{ voor } x > 25 & Een maximum. \\
A(25) = 25 \cdot 50 = 1250 & In m$^2$: de weide is $25$ m breed en $50$ m lang.
\end{stappenvrij}}{%
\begin{tikzpicture}[scale=.5,font=\scriptsize]
  \fill[accent!10] (0,0) rectangle (5,2.5);
  \draw[accent,thick] (0,0) -- (0,2.5) -- (5,2.5) -- (5,0);
  \draw[grijs] plot[domain=-.8:5.8,samples=80,smooth] (\x,{.12*sin(360*\x/1.1)});
  \node[below] at (2.5,0) {rivier};
  \node[left] at (0,1.25) {$x$};
  \node[above] at (2.5,2.5) {$100 - 2x$};
\end{tikzpicture}}}

Teken je de grafiek van $A(x)$, dan zie je een bergparabool met nulwaarden $0$ en $50$ en de top in $x = 25$. Aan de randen van het domein is de oppervlakte $0$: een weide zonder breedte, of zonder lengte.

\begin{werkwijze}
Kies een veranderlijke en druk de grootheid die optimaal moet zijn uit als functie van die ene veranderlijke (gebruik de voorwaarde om andere grootheden te elimineren). Bepaal het domein in de context. Leid af, zoek de nulwaarden, en controleer met een tekentabel (of de randen van het domein) dat het een maximum of minimum is. Beantwoord de vraag in de context.
\end{werkwijze}

\begin{oefeningen}
\oef{h4-weide2}
\end{oefeningen}

\voorbeeld{}{Een cilindervormig blik moet $1$ liter ($1000$ cm$^3$) bevatten. Welke afmetingen vragen het minste metaal?\label{vb:h4-blik}}{%
\uitleg{Minimaliseer de totale oppervlakte (bodem, deksel en mantel). Druk de hoogte $h$ uit met de inhoud: $\pi r^2 h = 1000$.}
\begin{stappenbreed}
h = \frac{1000}{\pi r^2},\quad A(r) = 2\pi r^2 + 2\pi r h = 2\pi r^2 + \frac{2000}{r} & $r > 0$. \\
A'(r) = 4\pi r - \frac{2000}{r^2} = 0 \iff r^3 = \frac{500}{\pi} & \\
r = \sqrt[3]{\frac{500}{\pi}} \approx 5{,}42,\quad h = \frac{1000}{\pi r^2} \approx 10{,}84 & $A'$ wisselt van $-$ naar $+$: minimum.
\end{stappenbreed}
Het zuinigste blik is even hoog als breed: $h = 2r$. Want uit $r^3 = \frac{500}{\pi}$ volgt $h = \frac{1000}{\pi r^2} = \frac{2r^3}{r^2} = 2r$.}

\begin{oefeningen}
\oef{h4-optimaal}
\end{oefeningen}

%========================================================================================
\opdrachtenkop

\begin{oefeningen}
\needspace{16\baselineskip}
\oef{h4-o10}

\oef{h4-o2}

\needspace{10\baselineskip}
\oef{h4-o16}

\oef{h4-o12}

\oef{h4-o1}

\needspace{14\baselineskip}
\oef{h4-o15}

\needspace{14\baselineskip}
\oef{h4-o11}

\oef{h4-o3}

\oef{h4-o13}

\oef{h4-o4}

\oef{h4-o14}

\needspace{12\baselineskip}
\oef{h4-o23}

\oef{h4-o19}

\oef{h4-o17}

\oef{h4-o5}

\oef{h4-o20}

\needspace{14\baselineskip}
\oef{h4-o24}

\oef{h4-blik2}

\oef{h4-o21}

\oef{h4-o6}

\oef{h4-o18}

\needspace{10\baselineskip}
\oef{h4-o22}

\oef{h4-o7}

\oef{h4-o26}

\oef{h4-o25}

\oef{h4-o8}

\oef{h4-o27}

\oef{h4-o9}
\end{oefeningen}

%----------------------------------------------------------------------------------------
% KOPPEL: verbanden met andere leerstof
%----------------------------------------------------------------------------------------
\needspace{20\baselineskip}
\opdrachtsectie{Koppel}

\lkr{Statistiek en Exponentiële en logaritmische functies: kijk na wat al gezien is. Het pakket legt de volgorde niet vast.}

\begin{oefeningen}
\oef{h4-koppel-1}

\oef{h4-koppel-3}
\end{oefeningen}

%----------------------------------------------------------------------------------------
\opdrachtsectie{Test jezelf}

\begin{oefeningen}
\oef{h4-t1}

\oef{h4-t2}

\oef{h4-t3}

\oef{h4-t4}

\oef{h4-t5}

\oef{h4-t6}

\oef{h4-t7}
\end{oefeningen}

\herhaalschema{%
  \fv{oef:h4-t1}{vb:h4-extrema}{oef:h4-tekentabel}%
  \fv{oef:h4-t2}{vb:h4-extrema2}{oef:h4-extremaexp}%
  \fv{oef:h4-t3}{vb:h4-buigpunt}{oef:h4-buigpunt}%
  \fv{oef:h4-t4}{vb:h4-weide}{oef:h4-weide2}%
  \fzv{oef:h4-t5}{oef:h4-o2}%
  \fv{oef:h4-t6}{vb:h4-buigpunt}{oef:h4-o10}%
  \fv{oef:h4-t7}{vb:h4-interval}{oef:h4-o22}}

\ifzelfstudie
\section*{Maak je eigen samenvatting}
Sluit dit boek en schrijf in je schrift: het verband tussen het teken van $f'$ en stijgen of dalen, wanneer een nulwaarde van $f'$ een maximum of minimum is, wat $f''$ zegt over hol en bol en wat een buigpunt is, en de werkwijze voor optimaliseringsproblemen. Vergelijk daarna met de kaders.
\fi
